\section*{Anexo II. Distribución de actuaciones del Equipo Pericial}
\addcontentsline{toc}{section}{EVIDENCIAS}
A efectos de trazabilidad interna y transparencia metodológica, se detalla a continuación la distribución de responsabilidades técnicas, custodia de soportes y análisis forenses ejecutados de forma colegiada por cada uno de los integrantes del \textbf{Equipo Pericial \GrupoPracticas} en la elaboración del presente dictamen:

\vspace{0.3cm}
\renewcommand{\arraystretch}{1.35}
\begin{tabularx}{\textwidth}{@{} >{\raggedright\arraybackslash}p{3.9cm} >{\raggedright\arraybackslash}p{3.5cm} X @{}}
	\toprule
	\textbf{\textcolor{AzulPericial}{Perito / Colegiado}} & 
	\textbf{\textcolor{AzulPericial}{Área de especialidad}} & 
	\textbf{\textcolor{AzulPericial}{Actuaciones y evidencias asignadas}} \\
	\midrule
	\textbf{\PeritoUno} \newline {\small Col. n.\textsuperscript{o} \NumColegiadoUno} & 
	Arquitectura de sistemas y soportes extraíbles & 
	\textbf{Escenario 3 (Dispositivo USB):} Custodia del soporte físico \texttt{General UDisk}, habilitación de bloqueo contra escritura (\texttt{WriteProtect = 1}), adquisición bit a bit en formato RAW (\texttt{usb\_evidencia.001}) con \textit{FTK Imager} y recuperación de \texttt{Confidencial.pdf} desde la \texttt{\$MFT} mediante \textit{Autopsy} (Sección 7). \\
	\midrule
	\textbf{\PeritoDos} \newline {\small Col. n.\textsuperscript{o} \NumColegiadoDos} & 
	Redes, protocolos SMTP y correo electrónico & 
	\textbf{Escenario 1 (Correo electrónico):} Custodia del disco virtual \texttt{Maquina\_Corrreo.vdi}, descifrado de volumen BitLocker en solo lectura (\textit{dislocker} en Kali Linux), extracción del mensaje MIME (\texttt{.eml}), análisis de \texttt{envios.log} y reconstrucción cronológica (Sección 5). \\
	\midrule
	\textbf{\PeritoTres} \newline {\small Col. n.\textsuperscript{o} \NumColegiadoTres} & 
	Sistemas de archivos NTFS y entornos en la nube & 
	\textbf{Escenario 2 (Dropbox):} Examen forense de la imagen \texttt{W11.dd}, inspección de la entrada MFT \texttt{188015} mediante \textit{Autopsy} y \textit{The Sleuth Kit} (\texttt{istat}), verificación del flujo alternativo \texttt{com.dropbox.attrs} y trazabilidad en \texttt{\$Recycle.Bin} (Sección 6). \\
	\midrule
	\textbf{\PeritoCuatro} \newline {\small Col. n.\textsuperscript{o} \NumColegiadoCuatro} & 
	Seguridad de la información y cadena de custodia & 
	\textbf{Infraestructura, metodología y consolidación:} Implementación y auditoría del servidor SMTP/Webmail de pruebas (\texttt{IFOR\_P2}), control de verificación criptográfica (SHA-256), adecuación normativa (ISO/IEC 27037 y UNE 71506) y síntesis de conclusiones (Secciones 3, 4, 8 y 9). \\
	\bottomrule
\end{tabularx}